\chapterhead{Introduction}

\sectionhead{Background of the Study}

Schools lend laboratory equipment to their students and staff, and keeping track
of what has gone out is a recognised management problem
\parencite{cepeda2025}. \textcite{cepeda2025} examined equipment lending at a
Philippine university and reported that handling the process by hand left the
institution weak in three areas at once: operational efficiency, data accuracy,
and accountability. Those three weaknesses are the reason such institutions move
to recorded, automated lending \parencite{cepeda2025}. The gap survives because
entry to a room and the record of what left it are governed by two separate
mechanisms that never verify one another \parencite{cepeda2025}.

Records maintained by hand also drift away from the physical stock they are
meant to describe \parencite{linuwih2025}. \textcite{linuwih2025} counted stock
against the written record in a warehouse and found that 52.85 percent of
product types did not match during a single stocktake. A written record that
disagrees with the shelf in more than half of cases cannot settle a question
about who holds what \parencite{linuwih2025}. Physical keys carry a separate
weakness of their own \parencite{alrawili2024}. \textcite{alrawili2024} observe
that anything a person carries, as opposed to something a person is, can be
handed to somebody else. A key is therefore able to be lent, copied, or lost,
after which the lock records only that somebody held a key rather than who that
person was \parencite{alrawili2024}.

Biometric identification answers that weakness directly, because a fingerprint
belongs to the person rather than being carried by them
\parencite{alrawili2024}. \textcite{alrawili2024} reviewed how such systems are
chosen and found that the choice is always a balance. A system has to be
secure, but it also has to be convenient and robust, or it will not survive
contact with the people who use it \parencite{alrawili2024}. Fingerprints sit at
a workable point on that balance, being inexpensive to read and quick enough to
complete a scan in roughly one second on low-cost hardware
\parencite{alsayaydeh2025}.

Fingerprint locks built on low-cost microcontrollers are now a well documented
application of this hardware \parencite{rahman2026}.
\textcite{alsayaydeh2025} built one using an ESP32, a small computer chip with
built-in Wi-Fi. Their lock read a finger in roughly one second and released the
door about five seconds later \parencite{alsayaydeh2025}. It admitted the wrong
person 1.32 percent of the time, which is a good result
\parencite{alsayaydeh2025}. It also turned away the right person 26.32 percent
of the time, so nearly one legitimate user in four had to present a finger again
\parencite{alsayaydeh2025}. That rejection figure is the convenience side of the
balance described by \textcite{alrawili2024}, on which a system that is not
convenient will not survive contact with its users.

A single reader also carries a limit that no improvement in sensor accuracy can
remove, since one reader verifies only the credential presented to it and
nothing about the person who may have entered alongside its user
\parencite{mohammed2025}. Researchers have answered this by verifying a person
more than once rather than by trying to perfect a single check
\parencite{mohammed2025}.
\textcite{mohammed2025} tested a system built this way
and measured a 68 percent drop in fraudulent access. The gain came from
requiring more than one independent confirmation before granting entry, not from
a more accurate sensor \parencite{mohammed2025}. Their framework still answered
in 235 milliseconds on average, so the extra check did not make the system slow
\parencite{mohammed2025}.

The record itself must also resist alteration, since a saved list is only
useful if nobody can revise it after an incident \parencite{koisser2023}.
\textcite{koisser2023} studied this problem across connected devices and
describe tampering as the central weakness of digital records, which is why
their scheme is built to make any modification detectable rather than merely
discouraged. Read together, these findings describe one design: verify a person
twice and bind the second check to the first \parencite{mohammed2025}, then hold
the result where it cannot be edited \parencite{koisser2023}.

\sectionhead{Statement of the Problem}

A school stockroom cannot prove who took its materials, because the
accountability its records are meant to provide is precisely what manual
handling fails to deliver \parencite{cepeda2025}. Two weaknesses produce this
failure. A single reader verifies only the credential presented to it, so a
person admitted alongside an authorized user is never evaluated and never
recorded, which is the weakness \textcite{mohammed2025} address by requiring
more than one confirmation. Written records then disagree with the physical
stock they describe in more than half of the product types examined
\parencite{linuwih2025}, and moving those records onto a computer does not
repair the problem, since a list that can still be revised afterwards carries
the same weakness as a paper one \parencite{koisser2023}. The institution is
therefore left unable to show either that its list of names is complete or that
nothing on it was altered later, which are the two properties an accountable
record requires \parencite{koisser2023, linuwih2025}.

Existing work addresses the parts of this problem separately, and the gap
lies between them \parencite{alsayaydeh2025, mohammed2025, koisser2023,
cepeda2025}. Fingerprint locks on low-cost microcontrollers guard a single
control point \parencite{alsayaydeh2025}. Requiring more than one confirmation
is known to cut fraudulent access, but that result was obtained by combining
modalities at one point on server-class hardware rather than by sequencing them
across two physical locations \parencite{mohammed2025}. Tamper-evident logging
is well developed, and recent work brings it within reach of resource-limited
devices \parencite{yagiz2026}, yet it is treated as a logging problem rather
than one bound to a physical access decision \parencite{koisser2023}. Per-device
credentials and replay resistance have likewise been demonstrated on ESP32-class
hardware in isolation \parencite{lkhluf2026}. Equipment borrowing has been moved
onto computers, but the written claim is not tied to a physical opening
\parencite{cepeda2025}, which is why records still drift from the shelf
\parencite{linuwih2025}. Across the work reviewed here
\parencite{alsayaydeh2025, mohammed2025, koisser2023, yagiz2026,
lkhluf2026, cepeda2025, linuwih2025}, no single reported system joins these
into one scheme: sequential, identity-bound verification across two control
points, running on low-cost microcontrollers, with an unalterable record
that binds what was taken to the opening that allowed it.

Specifically, this study seeks to answer the following questions:

\begin{enumerate}
  \item Can sequential, identity-bound verification across two control points be
        implemented on low-cost microcontrollers, given that the reduction in
        fraudulent access reported by \textcite{mohammed2025} was obtained on
        server-class hardware?
  \item What session lifetime binds the second check to the first without
        obstructing legitimate use, given the rejection rate reported for
        single-reader fingerprint locks \parencite{alsayaydeh2025} and the
        idle-timeout range recommended for high-value applications
        \parencite{owasp2026}?
  \item Can an access record be held unalterable by every role, including
        administrators, on a system of this scale, where tamper-evident logging
        has so far been treated separately from the access decision itself
        \parencite{koisser2023, yagiz2026}?
  \item Can a written record of borrowed items be bound to a physical opening,
        closing the gap left when borrowing is digitised without such a link
        \parencite{cepeda2025, linuwih2025}?
  \item Do the processor, memory, input/output, and communication resources of
        the ESP32 support this combination, given the limited on-chip memory
        identified as the platform's standing constraint \parencite{rahman2026}?
  \item What residual weaknesses remain in a design that delegates matching to
        the sensor, given that low-cost fingerprint modules have been shown to
        be spoofable \parencite{rad2024}?
\end{enumerate}

\sectionhead{Objectives of the Study}

The general objective of this study is to design and build an embedded system
that can prove who took materials from a school stockroom, and can store that
proof in a form that nobody is able to change afterwards.

The specific objectives below answer the six questions above, in the same
order, each stating what the study builds or measures to settle its
matching question.

\begin{enumerate}
  \item To implement sequential, identity-bound verification on two low-cost
        microcontroller readers, so that the cabinet opens only for the person
        the door admitted, and to demonstrate it on assembled hardware rather
        than on server-class equipment \parencite{mohammed2025}.
  \item To set and enforce a door session lifetime of 120 seconds, taken from
        the 2 to 5 minute idle timeout recommended for high-value applications
        \parencite{owasp2026}, and to close it automatically.
  \item To hold every scan, whether allowed or refused, in a record that offers
        no route to edit or delete it for any role, administrators included
        \parencite{koisser2023}.
  \item To bind a written list of borrowed items to a physical opening by
        requiring a one-time code that only the opening cabinet displays
        \parencite{cepeda2025}.
  \item To document how the processor, memory, input/output, and communication
        resources of the ESP32 support this combination, and to measure what the
        finished program costs against the limited on-chip memory that
        \textcite{rahman2026} identify as the platform's standing constraint.
  \item To state the residual weaknesses of the design, including those
        inherited by delegating fingerprint matching to a sensor module of a
        class shown to be spoofable \parencite{rad2024}.
\end{enumerate}

\sectionhead{Significance of the Study}

This study is useful to the groups described below.

\textbf{Students.} Students who borrow equipment gain a fair and complete
record. Because every scan is saved with a name and a time, a student cannot
easily be blamed for something that somebody else removed. The record protects
the borrower just as much as it protects the school, and it removes the argument
that so often follows a missing item. Students also gain a faster way of
borrowing, since there is no notebook to find and no queue behind the person
writing in it.

\textbf{Teachers and laboratory staff.} Staff no longer need to read other
people's handwriting or chase students for entries that were never made. The
list builds itself as people use the room, and it can be searched in seconds
rather than flipped through page by page. Staff can also see at a glance whether
a reader has stopped working, instead of discovering it only when something goes
wrong.

\textbf{Stockroom custodians and administrators.} Administrators gain something
a notebook can never provide, which is proof. When an item goes missing they can
see exactly who opened the cabinet, at what time, and what was logged against
that opening. This speaks directly to the record inaccuracy measured by
\textcite{linuwih2025}, where more than half of the product types checked
disagreed with their written record. Administrators can also see refused
attempts, which reveals whether people are trying to enter when they should not
be. That second view turns a vague worry about losses into a number that can be
watched over time.

\textbf{The school.} Laboratory equipment is expensive and school budgets are
limited. Better records mean fewer unexplained losses, and they support better
decisions about what needs replacing and how much to buy. Over several terms,
the record also shows which items are used heavily and which sit untouched.

\textbf{Future researchers.} This study documents what worked and what did not,
including a fault discovered during testing and the measurements still to be
taken. Other students can reuse the design, extend it to more rooms, or take the
honest account of its limits as a starting point and avoid repeating the same
mistakes.

\sectionhead{Scope and Limitations}

This study covers one door and one storage cabinet inside a single school
stockroom, served by two reader units and one server. A server is the computer
that stores all the records and makes the decisions about who may enter. The
system handles signing up a fingerprint, the two scans, the borrowing record,
the saved history, and text message alerts. It also shows different information
to students, teachers, and administrators, so that a student sees only their own
records while an administrator sees everything.

The study carries the following limits, and every claim made later in this paper
should be read against them.

\begin{enumerate}
  \item \textbf{Performance was not measured in numbers.} Both readers were
        built, wired, flashed and exercised through the complete door-to-cabinet
        sequence, so the system is demonstrated to work. What the study does not
        provide is quantitative performance data: scan time, unlock time and
        error rates were not recorded over a controlled set of trials, so no
        such figure is claimed here.
  \item \textbf{The sensor does the matching.} Fingerprints are stored and
        compared inside the sensor itself, never on our chip. The system is
        therefore only as accurate as the sensor, and it has no way of telling a
        real finger from a fake one. \textcite{rad2024} showed that copies made
        from seaweed jelly can fool this kind of sensor.
  \item \textbf{People must sign up twice.} Once at the door and once at the
        cabinet, because each sensor keeps its own separate list of fingers.
  \item \textbf{The borrowed list is still typed by a person.} The system can
        prove that the cabinet opened and who opened it. It cannot check whether
        the list of items typed afterwards is complete or honest.
  \item \textbf{Text messages were not sent over a real network.} They were
        tested using a stand-in program on the computer rather than a real SIM
        card and a real phone line.
  \item \textbf{One room only.} The design was not tested across several rooms
        or with several cabinets running at the same time.
\end{enumerate}

\sectionhead{Definition of Terms}

The terms below are listed in alphabetical order and are defined as they are
used in this study.

\begin{description}
  \item[Audit trail] The saved list of every scan and every decision the system
        has made.
  \item[Embedded system] A small computer built into a device to do one
        particular job.
  \item[ESP32] The chip used in this project. It is a small computer with
        built-in Wi-Fi.
  \item[Fingerprint sensor] The part that reads a finger and compares it against
        a list stored inside itself.
  \item[Memory] Space inside the chip that holds the program and the data it is
        working on.
  \item[Relay] An electric switch that lets a small chip control a much stronger
        power supply.
  \item[Server] The main computer that keeps every record and decides who may
        enter.
  \item[Session] A short time window, 120 seconds in this study, that opens when
        somebody scans at the door.
  \item[Solenoid lock] An electric lock that opens while power is sent to it and
        locks again when the power stops.
  \item[Tailgating] Walking into a room behind somebody who has just opened the
        door, without scanning a finger.
\end{description}
